\documentclass[11pt,a4paper]{article}

\usepackage[brazil]{babel}
\usepackage[T1]{fontenc}
\usepackage{lmodern}
\usepackage{microtype}
\usepackage[margin=2.2cm]{geometry}
\usepackage{amsmath}
\usepackage{booktabs}
\usepackage{tabularx}
\usepackage{enumitem}
\usepackage{algorithm}
\usepackage{algpseudocode}
\usepackage[hidelinks]{hyperref}
\usepackage{xcolor}

\definecolor{destaque}{HTML}{1F4E79}
\setlist{nosep, leftmargin=*}
\setlength{\parskip}{0.4em}
\setlength{\parindent}{0pt}

% Frase de abertura de cada seção: conecta com a anterior.
\newcommand{\gancho}[1]{\textcolor{destaque}{\textit{#1}}\par}

\floatname{algorithm}{Algoritmo}
\algrenewcommand\algorithmicrequire{\textbf{Entrada:}}
\algrenewcommand\algorithmicensure{\textbf{Saída:}}
\algrenewcommand\algorithmicfor{\textbf{para cada}}
\algrenewcommand\algorithmicdo{\textbf{faça}}
\algrenewcommand\algorithmicif{\textbf{se}}
\algrenewcommand\algorithmicthen{\textbf{então}}
\algrenewcommand\algorithmicend{\textbf{fim}}
\algrenewcommand\algorithmicreturn{\textbf{devolva}}
\algrenewcommand\algorithmicfunction{\textbf{função}}

\title{\textbf{NPCs para o jogo da velha}\\[0.3em]
  \large Minimax e Monte Carlo Tree Search}
\author{Caio Trindade e Luisa Ferreira\\
  \small Inteligência Artificial para Jogos I --- Busca Adversarial}
\date{Outubro de 2026}

\begin{document}
\maketitle

\section{O jogo como problema de busca}

\gancho{Para um NPC decidir uma jogada, o jogo precisa virar algo que um
algoritmo consiga percorrer.}

Seguimos a formulação de jogos do livro \emph{Artificial Intelligence: A Modern
Approach} (Russell e Norvig, cap.~5) e a implementação de referência dele, o
\texttt{games4e.py} do \href{https://github.com/aimacode/aima-python}{aima-python},
traduzido para C++. Um jogo é definido por:

\begin{itemize}
  \item \textbf{estado inicial} $s_0$: tabuleiro vazio, X começa;
  \item $\textsc{ToMove}(s)$: de quem é a vez;
  \item $\textsc{Actions}(s)$: casas livres;
  \item $\textsc{Result}(s, a)$: o tabuleiro depois da jogada $a$;
  \item $\textsc{TerminalTest}(s)$: alguém fez três em linha ou o tabuleiro encheu;
  \item $\textsc{Utility}(s, p)$: $+1$ se $p$ venceu, $-1$ se perdeu, $0$ no empate.
\end{itemize}

Essa é a interface \texttt{Game}, comum aos dois NPCs. Um jogador é qualquer
função que recebe o jogo e o estado e devolve uma jogada; por isso o Minimax e
o MCTS podem jogar entre si ou contra uma pessoa sem conhecer um ao outro.

O jogo da velha é pequeno: há $9! = 362\,880$ sequências de jogadas no máximo e
$255\,168$ partidas possíveis. A árvore inteira cabe numa busca completa, o que
torna o Minimax viável sem nenhuma aproximação.

\section{Minimax}

\gancho{Se o adversário também joga bem, a melhor jogada é a que garante o
melhor pior caso.}

Os dois jogadores são \textsc{Max} (quem está decidindo) e \textsc{Min} (o
adversário). O valor minimax de um estado é a utilidade que \textsc{Max}
consegue garantir a partir dele, supondo que os dois jogam de forma ótima:
\[
\textsc{Minimax}(s) =
\begin{cases}
\textsc{Utility}(s, \textsc{Max}) & \text{se } s \text{ é terminal}\\
\max_{a} \textsc{Minimax}(\textsc{Result}(s, a)) & \text{se é a vez de } \textsc{Max}\\
\min_{a} \textsc{Minimax}(\textsc{Result}(s, a)) & \text{se é a vez de } \textsc{Min}
\end{cases}
\]

A decisão (\texttt{minimaxDecision}) escolhe a ação cujo estado resultante tem
o maior valor. O cálculo é uma busca em profundidade que desce até os estados
terminais e sobe os valores alternando máximo e mínimo. Com fator de
ramificação $b$ e profundidade $m$, o custo é $O(b^m)$ em tempo e $O(bm)$ em
memória; no jogo da velha, $b \le 9$ e $m \le 9$.

\section{Poda alfa-beta}

\gancho{Muitos ramos não podem mudar a decisão; não é preciso visitá-los.}

A poda alfa-beta (\texttt{alphabetaSearch}) chega ao mesmo valor do Minimax
visitando menos nós. Durante a busca, ela guarda dois limites:
$\alpha$, o melhor valor já garantido para \textsc{Max} no caminho até ali, e
$\beta$, o melhor para \textsc{Min}. Quando um nó de \textsc{Max} encontra
um valor $v \ge \beta$, \textsc{Min} nunca deixaria o jogo chegar ali, e os
irmãos restantes são descartados; o caso de \textsc{Min} com $v \le \alpha$ é
simétrico.

\begin{algorithm}[h]
\caption{Alfa-beta com corte (\texttt{alphabetaCutoffSearch}), lado de \textsc{Max}}
\begin{algorithmic}[1]
\Function{MaxValue}{$s, \alpha, \beta, \mathit{prof}$}
  \If{$\textsc{CutoffTest}(s, \mathit{prof})$} \Return $\textsc{Eval}(s)$ \EndIf
  \State $v \gets -\infty$
  \For{$a \in \textsc{Actions}(s)$}
    \State $v \gets \max(v, \textsc{MinValue}(\textsc{Result}(s, a), \alpha, \beta, \mathit{prof}+1))$
    \If{$v \ge \beta$} \Return $v$ \Comment{poda}\EndIf
    \State $\alpha \gets \max(\alpha, v)$
  \EndFor
  \State \Return $v$
\EndFunction
\end{algorithmic}
\end{algorithm}

\textsc{MinValue} é igual trocando $\max$ por $\min$, $-\infty$ por $+\infty$
e a poda por $v \le \alpha$. Na raiz, a jogada escolhida é a de maior
\textsc{MinValue}; em caso de empate, fica a primeira na ordem de
\textsc{Actions}, da casa $(1,1)$ à $(3,3)$.

\section{Dificuldade: corte por profundidade}

\gancho{Um Minimax completo nunca perde. Para ter níveis, ele precisa enxergar
menos.}

A versão com corte para de descer quando a profundidade passa de $d$
($\textsc{CutoffTest}$: $\mathit{prof} > d$ ou estado terminal). Nesse ponto,
o estado vale $\textsc{Eval}(s)$, que por padrão é a própria utilidade: $\pm1$
se o jogo acabou e $0$ caso contrário. Ou seja, o que fica além do horizonte é
tratado como empate. As jogadas da raiz são avaliadas com profundidade 1:

\begin{center}
\begin{tabular}{@{}lll@{}}
\toprule
Nível & Enxerga & Consequência \\
\midrule
$d = 0$ & só a própria jogada & ganha quando pode, mas não bloqueia \\
$d = 1$ & a própria jogada e a resposta & bloqueia vitórias imediatas \\
$d = 2$ & mais uma jogada própria & ainda não prevê garfos (duas ameaças) \\
completo & até o fim do jogo & joga de forma ótima, nunca perde \\
\bottomrule
\end{tabular}
\end{center}

Escolhemos a profundidade, e não uma probabilidade de jogada aleatória, porque
ela é o equivalente ao número de iterações do MCTS: nos dois casos, a
dificuldade é quanto o agente \emph{pensa} antes de jogar, e o confronto entre
eles compara o uso desse orçamento.

\section{Implementação}

\gancho{A tradução do aima-python manteve a lógica e mudou só os nomes para o
estilo do repositório.}

\begin{center}
\begin{tabularx}{\textwidth}{@{}l>{\raggedright\arraybackslash}X@{}}
\toprule
Arquivo & Conteúdo (nome no aima-python) \\
\midrule
\texttt{Game.h/.cpp} & \texttt{GameState}, \texttt{Game} e \texttt{playGame} (\texttt{play\_game}) \\
\texttt{TicTacToe.h/.cpp} & \texttt{TicTacToe} \\
\texttt{Minimax.h/.cpp} & \texttt{minimaxDecision}, \texttt{alphabetaSearch} e
  \texttt{alphabetaCutoffSearch} (\texttt{minimax\_decision},
  \texttt{alphabeta\_search}, \texttt{alphabeta\_cutoff\_search}) \\
\texttt{Players.h/.cpp} & \texttt{queryPlayer}, \texttt{randomPlayer} e
  \texttt{alphabetaPlayer} (\texttt{query\_player}, \texttt{random\_player},
  \texttt{alphabeta\_player}) \\
\bottomrule
\end{tabularx}
\end{center}

A única mudança de comportamento é no \texttt{queryPlayer}: ele pede a jogada
de novo quando ela é inválida. No original, uma jogada inválida não altera o
estado, mas passa a vez, e o próximo jogador joga com a peça errada.

Os jogadores são escolhidos pela linha de comando, X primeiro:
\texttt{./build/tictactoe minimax:1 aleatorio}. As opções são \texttt{humano},
\texttt{aleatorio}, \texttt{minimax} (completo) e \texttt{minimax:<d>}.

\section{Resultados do Minimax}

\gancho{Os níveis precisam ser diferentes na prática, não só na teoria.}

Cada nível jogou 200 partidas contra o jogador aleatório como X (primeiro a
jogar) e 200 como O.

\begin{center}
\begin{tabular}{@{}lrrrrrr@{}}
\toprule
 & \multicolumn{3}{c}{Como X} & \multicolumn{3}{c}{Como O} \\
\cmidrule(lr){2-4}\cmidrule(l){5-7}
Nível & Vitórias & Empates & Derrotas & Vitórias & Empates & Derrotas \\
\midrule
$d = 0$  & 176 & 2 & 22 & 121 & 2  & 77 \\
$d = 1$  & 192 & 5 & 3  & 149 & 31 & 20 \\
$d = 2$  & 189 & 6 & 5  & 152 & 29 & 19 \\
completo & 199 & 1 & 0  & 166 & 34 & 0 \\
\bottomrule
\end{tabular}
\end{center}

O Minimax completo nunca perdeu. A maior diferença está entre $d = 0$ e
$d = 1$: sem ver a resposta do adversário, o NPC não bloqueia e perde 77
partidas como O; vendo-a, perde 20. Entre $d = 1$ e $d = 2$ quase não há
diferença: o lance extra mostra as próprias vitórias um lance adiante, mas
ainda não chega aos garfos do adversário, que só aparecem olhando mais longe.
Jogar como O é mais difícil para todos os níveis, já que X começa.

Contra o Minimax completo, $d = 1$ e $d = 2$ empatam jogando como X, e $d = 0$
perde; jogando como O, todos os níveis limitados perdem. Como os dois lados são
determinísticos, cada confronto entre níveis do Minimax é sempre a mesma
partida.

% ---------------------------------------------------------------------------
% A completar: Monte Carlo Tree Search
% - seleção (UCB), expansão, simulação e retropropagação
% - dificuldade pelo número de iterações N
% - resultados do MCTS contra o jogador aleatório

\section{Monte Carlo Tree Search}

\gancho{}

% A completar: Minimax contra MCTS
% - partidas entre os níveis de profundidade d e de iterações N,
%   alternando quem começa
% - comparação e conclusão

\section{Minimax contra MCTS}

\gancho{}

\end{document}
